\cotefile{156-4}
\scanpage{82}\mspage{81}

\diagram{\begin{tikzcd}[ampersand replacement=\&,column sep=0.6em,row sep=0.15em]
\deleted{X\ll} \& \deleted{G} \& \deleted{\leq} \& \deleted{F}\\
\& \rotatebox{90}{$\deleted{\leq}$} \& \& \rotatebox{90}{$\deleted{\leq}$}\\
\& \deleted{G_X} \& \deleted{\leq} \& \deleted{X}
\end{tikzcd}}

\deleted{Soit $G_X$ la sous figure de $G$ image inverse de \deleted{\unsure{la|les}} sous figure $X$ de $F$. Par la prop. (p.~\unsure{76|70}) $X$ \unsure{raffine|raffinent} \unsure{à} : la fois la sous-figure $X$ de $F$, et la \unsure{raffine} $G_X$ \unsure{donc|dans} $X^{\unsure{2|<}}\ll G_X <$}

\deleted{Soit} Supp.\deleted{\unsure{sons|ons}} $X\leq F$, \unsure{i.e.|Id e.} $X$ ss-fig. de $F$, considérons ses images inverses \inserted{$\bar F_X,\deleted{\unsure{\bar G|G}}G_X,\bar G_X$} dans \deleted{\unsure{les}} $\bar F,G,\bar G$, \inserted{ce sont des sous-figures de $\bar F,G,\bar G$ \deleted{\unsure{resp.}}.} Donc on trouve \deleted{\unsure{(par la prop.}}
\[
\deleted{X\ll\bar G_X}
\]

\deleted{et on est ramené au cas particulier : $H$ ($=\bar G$) figure, $H'$ ($=\bar G_X$) sous-figure, $X$ multistrate de \deleted{$H$ (\unsure{i.e.|=} $X\leq H'$) \unsure{supposons}.}}

par la propriété universelle \unsure{du|de la} \unsure{passage}
\diagram{\begin{tikzcd}[ampersand replacement=\&,column sep=0.55em,row sep=0.15em]
G_X \& \leq \& X\\
\rotatebox{90}{$\unsure{\preceq|\ll}$} \& \& \rotatebox{90}{$\preceq$}\\
\bar G_X \& \leq \& \unsure{F|P}_X\\
\rotatebox{90}{$\leq$}\\
X
\end{tikzcd}}
On est ramené \unsure{:|à} prouver $X\leq G_X$

\deleted{D'où} \unsure{Com|Comm.} \unsure{par|ça} admettre
\deleted{$\unsure{X|G}\ \unsure{X|\ll}\ \unsure{\preceq|\leq}\ G$}
\margin{\deleted{$(\unsure{\leq|\ll}F_X\unsure{\leq|<})$}}

% A short vertical rule stands to the left of a) and b).
\noindent\llap{\smash{\rule[-2.9\baselineskip]{.35pt}{3.4\baselineskip}}\hspace{.6em}}\deleted{\uline{C8} a) Soit \deleted{$X$ une multistrate \deleted{\unsure{$X$|de}}, $\Phi$ \unsure{une}} $X$ une multistrate, $\Phi$ une subdivision de $X$, si $X\leq\Phi$ (i.e. $X$ est une multistrate de $\Phi$) \unsure{on} a $X=\Phi$.}

\deleted{b) Si $X\preceq\Psi$, $X$ multistrate, $\Psi$ figure, alors $X=\Psi$}

\deleted{Donc dans le cas d'espèce on trouve}
\[
\deleted{X\ \deleted{\unsure{\leq|\preceq}}\ \leq\bar G_X\leq F_X=X}
\qquad\text{\deleted{donc}}\quad\deleted{X=\bar G_X}\quad\text{\deleted{, d'où}}
\]
\[
\deleted{X\preceq G_X\leq X}
\]

\scanpage{83}\mspage{82}

\noindent\llap{\smash{\rule[-4.9\baselineskip]{.35pt}{5.6\baselineskip}\hspace{.20em}\rule[-4.9\baselineskip]{.35pt}{5.6\baselineskip}}\hspace{.6em}}\uline{C 8}\textsubscript{(0)} Considérons les situations
% slanted
% Two parallel vertical strokes run down the left of C8.
\margin{\uline{provisoire !}}
\[
\begin{aligned}
\text{(a)}\qquad &X\leq\Phi\preceq X & \Phi&\in\unsure{\mathcal F|\mathscr F},\ X\in\unsure{\mathcal M|\mathscr M}\\
\text{(b)}\qquad &X\preceq\Psi\leq X & \deleted{\unsure{F|\Phi}}\Psi&\in\mathcal F,\ X\in\mathcal M
\end{aligned}
\]
\deleted{\unsure{C|A}}Alors on a $X=\Phi\quad\text{et}\quad X=\Psi$

\unsure{Achevons} la dém. Appliquant (a) à
\[
X\leq F_X\preceq X
\]
on trouve $X=F_X$, d'où, comme
\[
X\leq\bar G_X\leq F_X=X\quad,
\]
$X=\bar G_X$, d'où
\[
\deleted{\unsure{X=|X\leq}}\bar G_X=X\preceq G_X\leq X
\]
et appliquant (b), on trouve $\deleted{\unsure{G|X}}X=G_X$ et par suite $X\leq G_X$, qed.

\uline{Remarque} Bien sûr, C8 est nécessaire si on veut que $\ll$ soit une relation d'ordre. Et c'est aussi \unsure{suffisant}\unsure{,:|;:|;}

\uline{Cor} (de C8 \inserted{via prop. p\deleted{\unsure{g|.}}~\unsure{76|70}}) Si $F\ll G$ et $G\ll F$, on a $F=G$

\uline{Dém} \deleted{Soit $X\in\widetilde F$ i.e. $X\in\mathcal M$, $X\leq F$} On va prouver $\unsure{F|\bar F}\leq G$\ ; \deleted{\unsure{(on}} on aura par symétrie $G\leq F$, donc $F=G$.
\deleted{Il faut prouver \deleted{\unsure{pour|que}} pour $X\in\mathcal M$, $(X\unsure{\leq|<}F)\ \unsure{-|\Rightarrow|=} X\leq G$.}

\diagram{\begin{tikzcd}[ampersand replacement=\&,column sep=0.45em,row sep=0.1em]
\deleted{F} \& \deleted{\ll} \& \deleted{G} \& \deleted{\ll} \& \deleted{F}\\
\rotatebox{90}{$\deleted{\leq}$} \& \& \rotatebox{90}{$\deleted{\leq}$} \& \& \rotatebox{90}{$\deleted{\leq}$}\\
\deleted{F_X} \& \deleted{\ll} \& \deleted{G_X} \& \deleted{\ll} \& \deleted{X}\\
\rotatebox{90}{$\deleted{\leq}$}\\
\deleted{X}
\end{tikzcd}}

\scanpage{84}\mspage{83}

On a par définition de $\ll$
\diagram{\begin{tikzcd}[ampersand replacement=\&,column sep=0.65em,row sep=0.15em]
F \& \leq \& G' \& \preceq \& G\\
\& \& \& \& \rotatebox{90}{$\geq$}\\
\& \& \& \& F'\\
\& \& \& \& \rotatebox{90}{$\succeq$}\\
\& \& \& \& \deleted{G}F
\end{tikzcd}}

on complète par C\deleted{\unsure{6|8}}7
\diagram{\begin{tikzcd}[ampersand replacement=\&,column sep=0.6em,row sep=0.15em]
F \& \leq \& G' \& \preceq \& G\\
\& \& \rotatebox{90}{$\deleted{\unsure{\leq|\preceq}}\geq$} \& \& \rotatebox{90}{$\geq$}\\
\& \& H \& \preceq \& F'\\
\& \& \& \& \rotatebox{90}{$\succeq$}\\
\& \& \& \& \deleted{G}F
\end{tikzcd}}

\deleted{\unsure{d|S}} donc
% Both manuscript labels stand at the left of their equations.
\begingroup
\makeatletter
\tagsleft@true
\let\veqno\@@leqno
\makeatother
\[
F\leq H\preceq F \tag*{$(*)$}
\]
On va en conclure $F=H$ (\unsure{généralise} C8 (a)) d'où $G'=F$, $F'=F$, d'où
% slanted
\margin{Mais ce sont \uuline{les \unsure{\={m}s|mes|mêmes}} cas, en échangeant les \unsure{rôles|rôle} de $F$ et $H$ on trouve $H\preceq F\leq H$ !}
\[
F\preceq G\leq F \tag*{$(**)$}
\]
\endgroup
on va en conclure $F=G$ (\unsure{généralise} C8 (b)).

Prouvons d'abord ce dernier \unsure{pt|point} : il \deleted{faut} suffit de prouver $F\leq G$, i.e. \deleted{\unsure{Si|soit}} $X\in\unsure{\mathcal M|\mathscr M}$, $X\leq F\Longrightarrow X\deleted{\unsure{<|\leq}}\leq G$, c'est là\unsure{.|:} \unsure{ess.|cor.|ex.} le corollaire p.~\unsure{82|81}.

\deleted{Pour \unsure{traiter|traiter le} $(\unsure{*|**})$}

On a essayé ici de se débrouiller aussi loin que possible, sans \unsure{critères|critère} et raisonnements de disjonction, genre C5.

\scanpage{85}\mspage{84}

\uline{Remarque} Soi\deleted{\unsure{en|n}}t $\unsure{\mathcal F|\mathscr F}$ \deleted{\unsure{un|ens.}} muni de deux \unsure{relations|relation} d'ordre $\leq$, $\preceq$, satisfaisant les conditions suivantes

a) Pour
\diagram{\begin{tikzcd}[ampersand replacement=\&,column sep=0.55em,row sep=0.25em]
F \& \deleted{\unsure{\leq|\preceq}}\geq \& G\\
\rotatebox{90}{$\preceq$}\\
F'
\end{tikzcd}}
on peut compléter de façon unique en
\diagram{\begin{tikzcd}[ampersand replacement=\&,column sep=0.65em,row sep=0.2em]
\& F \& \geq \& G\\
(*) \& \rotatebox{90}{$\preceq$} \& \& \rotatebox{90}{$\preceq$}\\
\& \deleted{\unsure{G|F}}F' \& \geq \& G'
\end{tikzcd}}

b) Pour
\diagram{\begin{tikzcd}[ampersand replacement=\&,column sep=0.55em,row sep=0.2em]
\deleted{F}G \& \leq \& F\\
\rotatebox{90}{$\preceq$}\\
G'
\end{tikzcd}}
, on peut compléter (de façon pas \unsure{néc\textsuperscript{t}|néc'.} unique) en
\diagram{\begin{tikzcd}[ampersand replacement=\&,column sep=0.65em,row sep=0.2em]
G \& \leq \& F\\
\rotatebox{90}{$\preceq$} \& \& \rotatebox{90}{$\preceq$}\\
G' \& \leq \& F'
\end{tikzcd}}

c) Si
\[
\begin{aligned}
&F\leq G\preceq F &&\text{(ou ce qui revient au }\unsure{=|\bar m}\\
&G\preceq F\leq\deleted{\unsure{F|I}}G &&),
\end{aligned}
\]
alors $F=G$.

\uline{Alors} :
\deleted{(i) Posant $F\ll\ \overset{\mathrm{def}}{\Longleftrightarrow}\ \exists$}
(i) \unsure{Posant} $\deleted{F}G\ll\deleted{G}F\ \overset{\mathrm{def}}{\Longleftrightarrow}\ \exists\deleted{G}F'$ \unsure{avec|avc} $G\leq F'\preceq F$
la relation $\deleted{\unsure{\leq|<}}\ll$ est une relation d'ordre (et de préordre, \unsure{si on|si} ne suppose que a) b) )

\scanpage{86}\mspage{85}

(ii) Dans le cas a), le carré $(*)$ est cartésien pour $\ll$, i.e.
\[
G'=\underset{\ll}{\mathrm{Inf}}(F',G)
\]
% slanted
\margin{Il manque \deleted{\unsure{peut-être}} \inserted{sans doute} des hypothèses ?? P.ex. que les $\mathrm{Inf}$ finis existent pour $\leq$, et\ldots\ Il \uline{faudrait tirer} au \uline{clair la situation}}

\deleted{\unsure{(iii)|(ii)}} (iii) Si $\deleted{\unsure{G|G'}}$
\diagram{\begin{tikzcd}[ampersand replacement=\&,column sep=0.7em,row sep=0.2em]
\deleted{\unsure{X|G}}F \& \deleted{\unsure{\leq|\preceq}}\geq \& G\\
\rotatebox{90}{$\leq$} \& \& \rotatebox{90}{$\deleted{\unsure{\preceq|\leq}}\leq$}\\
F' \& \geq \& G'
\end{tikzcd}}
est cartésien pour \rotatebox{90}{$\leq$}, il est cartésien pour $\deleted{\unsure{\leq|\preceq}}\ll$.

Dans les axiomes des \unsure{cartésies|cartésiens|cartiers}, il faudrait séparer

\tikz[baseline=(n.base)]\node[draw,circle,inner sep=0.7pt](n){$\unsure{\alpha|*|K}$}; un paquet concernant $\leq$ tout seul
\inserted{(C1 \unsure{à|:} C4)}

\tikz[baseline=(n.base)]\node[draw,circle,inner sep=0.7pt](n){$\unsure{\beta|8|*}$}; un paquet \unsure{autour|entier} des propriétés ``formelles'' \unsure{reliant} $\leq$, $\preceq$ (\unsure{ess\textsuperscript{t}.|ess!} celles qu'on vient d'esquisser) sans questions de support ni de disjonction, et en utilisant le minimum sur $\leq$ tout seul.

Conditions qui assurent, quand les \unsure{sup|Sup} majorés existent pour $\leq$, que ceux ci commutent à l'induction sur une ``subdivision.

Au besoin, si on \unsure{ne|en} veut se donner la peine de \deleted{\unsure{d|l}}dégager les axiomes \unsure{minimum|minimums}, \unsure{donner|donnant} le paquet des propriétés formelles ``évidentes'' du pt de vue signification géom.

\scanpage{87}\mspage{86}

\tikz[baseline=(n.base)]\node[draw,circle,inner sep=0.7pt](n){$\unsure{*|\gamma}$}; Axiome de \uline{recollement} des subdivisions

\tikz[baseline=(n.base)]\node[draw,circle,inner sep=0.7pt](n){$*$}; \deleted{\unsure{A|U}}Un paquet d'axiomes de disjonction

\deleted{\unsure{C6 C8|C8}}\qquad -\quad -\quad -

\uline{C8}\quad \uline{Axiome} \uline{de recollement} \uline{des subdivisions}

% A vertical stroke sets off the following four manuscript lines.
\noindent\llap{\smash{\rule[-2.5\baselineskip]{.35pt}{3.1\baselineskip}}\hspace{.6em}}Soit $\unsure{\mathbb F|F}= \operatorname{Sup} F_i$, et pour tt $i$, soit $F'_i$ une subdivision de $F_i$, telle que pour $i,j$, on ait $F'_i\mid F_i\cap F_j=F'_j\mid F_i\cap F_j$, alors $F'=\operatorname{Sup}F'_i$ existe et est une subdivision de $F$.

\deleted{Soit} (Devrait suivre C6\unsure{)|\}}, \unsure{Cela|cela} remplace C6 (c)

\uline{Cor} \unsure{Soient|Soit} \deleted{$F$ et $\unsure{F_i|F_a}$, alors} \inserted{$F=\coprod F_i$, alors}
\[
\operatorname{Subdiv}(F)\simeq\prod_i\operatorname{Subdiv}(F_i)
\]

\uline{Axiomes de disjonction et de support}.

Pour trouver les axiomes sous forme aussi forte et \unsure{sensible|saisissable} que possible, c'est le moment de \unsure{réintroduire|reintroduire|introduire} les lieux.

\deleted{\deleted{\uline{Définition} On appelle} \inserted{Prop.} Soit $F\in\unsure{\mathcal F|\mathscr F}$, \inserted{$F\ne\varnothing_{\unsure{F|\mathcal F}}$,} Conditions équivalentes}

\deleted{a) $F'\preceq F\Longrightarrow F'=F$ \qquad i.e. $\deleted{\unsure{F|E}}F$ est minimal dans $\mathcal F$ pour $\preceq$}

\deleted{b) $F'\ll F$, \inserted{$F'\ne\varnothing$} $\Longrightarrow F'=F$ \qquad i.e. $F$ est minimal dans $\mathcal F$ \inserted{$\unsure{-|\setminus}\{\varnothing\}$} pour $\ll$}

\deleted{et ceci implique que $F\in\unsure{\mathcal M|\mathscr M}$ \deleted{\illegible{1}} et $F$ est minimal dans $\deleted{\unsure{\mathcal F|F}}\mathcal M$ (\unsure{ou en|ou dans} $\deleted{\unsure{\mathcal F|F}}\unsure{\mathcal M|\mathcal F}$) pour $\leq$,}

\scanpage{88}
\mspage{87}
\deleted{\unsure{87|8}}

\deleted{\unsure{Dém.|Dem.} $b)\Rightarrow a)$ trivial, prouvons $a\Rightarrow b$. Soit $F'\ll F$ i.e. $F'\leq\bar F\preceq F$. Par $a$, on a $\bar F=F$, donc $F'\leq F$. \deleted{Par $C7$, $\exists$} On a donc $F'\leq F$ \deleted{et $F'\unsure{\preceq|\leq}F$, \unsure{on veut|on} que ça implique $F'=F$.} \deleted{Mais si \unsure{l|le}} On va le prouver \unsure{tantôt|tantot}.}

\deleted{Ceci dit, comme $\ll$ \inserted{est} \unsure{impliqué|impliqués} \inserted{par} $\leq$, il en résulte que (si $F\neq\varnothing$) $F$ est \deleted{minimal dans $\unsure{\mathrm{Supp}|{\mathcal F}}\unsure{\{\varnothing\}|\{x\}|\setminus\{\varnothing\}}$ pour $\leq$ : Si $X'<$}}

% Single rule beside the definition.
\noindent\llap{\smash{\rule[-2.1\baselineskip]{.35pt}{2.8\baselineskip}}\hspace{.6em}}\uline{Définition} On appelle \uline{lieu} \inserted{de $\unsure{\mathcal E|\mathcal C}$} un élément de $\unsure{\mathcal F|\mathscr F}\setminus\{\varnothing_{\mathcal F}\}$ minimal pour $\ll$, i.e. une figure $\deleted{\unsure{X|x}}F$ telle que $\deleted{\unsure{X|x}}F\neq\varnothing_{\mathcal F}$, et $F'\ll F$, $F'\neq\varnothing_{\unsure{F|\mathcal F}}$ implique $F'=F$.

Ainsi, on voit qu'un lieu est néc. $\in\unsure{\mathcal M|\mathscr M}$, et c'est aussi un élément de $\mathcal M$ minimal pour $\leq$ (mais il n'est pas vrai, dans l'état actuel des choses, -- i.e. avec les axiomes qu'on \unsure{a} introduits \unsure{jusqu'à|jusqu'} présent) l'inverse soit vrai. (Ce sera vrai cependant pour les ``\unsure{cartésies|cartésiens} en boules''\ldots)

On désigne par $\unsure{\mathcal L|\mathscr L}$ l'ensemble des lieux.

\deleted{On} Pour \unsure{toute|tout} $F\in\mathcal F$, on appelle \uline{ombre de $F$} l'ens. des $x\in\mathcal L$ tels que $x\ll F$, notation $\lvert F\rvert$, et \uline{multiombre} l'ens
\scanpage{89}
\mspage{88}
des parties $\deleted{\unsure{\widehat F|\widehat H}}$ de $\unsure{\mathcal L|\mathscr L}$
\[
\widehat F=\bigl\{\,\lvert X\rvert\mid X\in\bar F\quad\text{i.e. }X\in\unsure{\mathcal M|\mathscr M},\,X\leq F\,\bigr\}
\]

Ainsi on a
\[
\lvert\overset{\deleted{\unsure{\wedge|{-}}}}{F}\rvert\subset\mathcal L,\qquad
\widehat F\subset\unsure{\mathcal P|P}(L)
\]

et on a
\[
\lvert F\rvert=\deleted{\unsure{\coprod|\mathcal L}}\bigcup_{X\in\bar F}\lvert X\rvert=\bigcup_{A\in\widehat F}A\,.
\]

Cette relation signifie que si $x\unsure{\in|+}\mathcal L$ est tel que $x\ll F$, $\exists\,X\in\bar F$ tel que $x\ll X$, \unsure{\S|(*)} 
% slanted
\margin{Mais pour \unsure{le|la} prouver, on a utilisé \unsure{axiome|l'axiome} \unsure{d'invention|d'induction|d'inversion}}
et pour ceci, voir prop. page \unsure{78|76}. Il y a \unsure{même|\={m}} un plus petit tel $X$ (\unsure{=|:} Cor à la prop).

\uline{Proposition} La multi-ombre $\unsure{\Phi|\mathcal P}=\widehat F$ d'une figure est une famille \uline{admissible} de parties de $\mathcal L$, i.e.
\[
\begin{aligned}
a)\quad &\forall\,\deleted{\unsure{a\in|x\in|A\in}}A\in\Phi,\ \text{posant }\ \unsure{\delta|\partial}A=\bigcup_{\substack{B\subset A\\B\neq A}}B\\
&\text{on a}\qquad \delta A\neq A\quad\text{\unsure{i.e|i.e.}}\quad A^{\unsure{\circ|o}}\overset{\mathrm{def}}{=}A\setminus\delta A\neq\varnothing
\end{aligned}
\]

\[
b)\quad \forall\,A,B\in\Phi,\ \text{on a}\qquad A\cap B=\bigcup_{\substack{C\in\Phi\\C\subset A\cap B}}C
\]
